% Generated draft: requires mathematical and source review.
\documentclass[11pt]{article}
\usepackage[margin=25mm]{geometry}
\usepackage{fontspec}
\setmainfont{Latin Modern Roman}
\usepackage{amsmath,amssymb,mathtools,unicode-math}
\setmathfont{Latin Modern Math}
\usepackage{enumitem,needspace,xurl,hyperref}
\setlength{\parindent}{0pt}
\setlength{\parskip}{5pt}
\setlength{\emergencystretch}{3em}
\newcommand{\N}{\mathbb{N}}
\newcommand{\Z}{\mathbb{Z}}
\newcommand{\Q}{\mathbb{Q}}
\newcommand{\R}{\mathbb{R}}
\newcommand{\C}{\mathbb{C}}
\newcommand{\F}{\mathbb{F}}
\newcommand{\A}{\mathbb{A}}
\newcommand{\PP}{\mathbb{P}}
\newcommand{\E}{\mathbb{E}}
\newcommand{\eps}{\varepsilon}
\newcommand{\dd}{\,\mathrm{d}}
\begin{document}
\section*{UnsolvedMath ID: 30003636}
\textbf{Source problem code:} catalogue entry
\subsection{Definitions and mathematical statement}
In the standard complex rank-one interpretation, let \(X\) be a smooth connected projective complex surface, let \(D\subset X\) be the support of a divisor, and put \(U=X\setminus D\). These choices provide a framework for the question; the supplied record does not explicitly specify the ground field, connectedness, or conventions concerning divisor multiplicities. For a base point \(u\in U\), define the character space and its cohomology jump loci by
\[
\mathbb T(U)=\operatorname{Hom}(\pi_1(U,u),\C^*),\qquad
\mathcal V^q_r(U)=\{\rho\in\mathbb T(U):\dim_{\C}H^q(U;\C_\rho)\ge r\},
\]
where \(q\ge0\), \(r\ge1\), and \(\C_\rho\) is the rank-one local system with monodromy \(\rho\). An irreducible component means an irreducible algebraic component of the specified jump locus.

\textbf{Source question.} Does the characterization of high-dimensional essential components of the characteristic variety for a union of lines in \(\mathbb P^2\) extend to more general divisors on smooth projective surfaces, or is that characterization specific to unions of lines?

\textbf{Unresolved formulation.} The supplied record does not state the line-arrangement characterization being referenced. It also does not identify the cohomological degree \(q\), the depth \(r\), the dimension threshold meant by ``high-dimensional,'' or the definition of ``essential.'' Consequently, the displayed jump loci explain a standard interpretation of characteristic varieties but do not identify the precise locus intended by the source. Neither the exact proposed equivalence nor its hypotheses for general divisors can be reconstructed from this record. The question is therefore retained as an extension problem requiring those specifications, rather than asserted as a fully quantified conjecture for every divisor and every surface.
\subsection{Short English statement}
The question asks whether a criterion describing high-dimensional essential components of characteristic varieties of plane line complements also applies to complements of more general divisors on smooth projective surfaces. The supplied record omits the criterion and several defining conventions, so a precise universal assertion remains unresolved.
\subsection{Sources}
\begin{itemize}
\item Supplied local UnsolvedMath registry record, permanent ID \(30003636\), title ``Characteristic Varieties of Divisors Beyond Line Arrangements,'' exact source problem code \texttt{catalogue entry}. The record belongs to the Oberwolfach Reports Open Problems collection and gives proposed year 2017; no individual proposer is supplied.
\item Supplied source citation: \emph{Mini-Workshop: Interactions between Low-dimensional Topology and Complex Algebraic Geometry} (2018), Oberwolfach Reports, pp.~2987--3033, DOI 10.4171/owr/2017/50. \url{https://doi.org/10.4171/owr/2017/50}. The background describes Oberwolfach Reports as documenting mathematical workshops and their open-problem discussions. The report citation is supplied provenance; its contents have not been independently consulted here.
\item Supplied literature reference: Cogolludo--Libgober, \url{https://arxiv.org/abs/1904.00969}. No paper title or theorem identifying the intended characterization is supplied.
\item Supplied literature reference: Oberwolfach Report 50/2017, \url{https://doi.org/10.4171/owr/2017/50}.
\item Literature assessment dated exactly \texttt{2026-08-21}: the supplied curation labels the problem ``partially solved'' and states that later work gives broader finiteness and free-quotient results for smooth quasiprojective settings, but does not establish a complete analogue of the line-arrangement criterion in every divisor and surface setting. This is supplied dated curation, not an independently verified assessment. The source itself requests independent verification of time-sensitive status and cited claims.
\end{itemize}
\end{document}
